\documentclass[12pt]{article}
\usepackage[utf8]{inputenc}
\usepackage[margin=0.75in]{geometry}
\usepackage{amsmath}
\usepackage{amssymb}
\usepackage{amsthm}
\usepackage{graphicx} % Required for inserting images
\usepackage{array}
\usepackage{tikz-cd}


\title{MAT1100 Homework 2}
\author{Aaron Avram}
\date{October 5 2026}

\begin{document}
\newtheorem{definition}{Definition}
\newtheorem{theorem}{Theorem}
\newtheorem{claim}{Claim}
\newtheorem{corollary}{Corollary}

\maketitle

\begin{align*}
    \textbf{Question 1}
\end{align*}

\begin{enumerate}
    \item[(a)]
    \begin{proof}
        Suppose $K, N \lhd G$ with $K \leq N$.

        First we show that
        \[ N/K \lhd G/K \]
        suppose $g \in G$ and $n \in N$ then
        \[ [g]_K[n]_K[g^{-1}]_K = [gng^{-1}]_K \in N/K \]
        where final part follows from the fact that $N \lhd G$ so
        $gng^{-1} \in N$.

        Next, to construct a isomorphism from
        \[ (G/K)/(N/K) \cong G/N \]
        it suffices to find a surjection
        \[ \varphi: G/K \to G/N \]
        whose kernel is $N/K$ and factor this map through the quotient by its kernel.

        Concretely, define
        \[ \varphi: [g]_K \to [g]_N \]
        first to see that $\varphi$ is well defined: suppose $g, h \in G$
        such that $[g]_K = [h]_K$ which means that $g^{-1}h \in K \leq N$.
        Therefore
        \[ [g]_N = [h]_N \]
        hence $\varphi$ is well defined.

        That $\varphi$ is surjective is clear, as if $g \in G$
        \[ \varphi([g]_K) = [g]_N \]
        by definition.

        Next, we show that $\ker \varphi = N/K$. It is clear
        that $N/K \subseteq \ker \varphi$. As if $n \in N$, $[n]_N = e$
        and $\varphi([n]_K) = e$. Next if $[g]_K \in \ker \varphi$
        then
        \[ [g]_N = e \]
        and hence $g \in N$ hence $[g]_K \in N/K$ as required.
        
        Now by the first isomorphism theorem
        \[ (G/K) / (N/K) \cong (G/N) \]
        with an explicit isomorphism given by
        \[ \hat{\varphi}: [[g]_K]_{N/K} \mapsto \varphi([g]_K) = [g]_N \]
    \end{proof}

    \item[(b)]
    \begin{proof}
        Suppose $K = \{ e, (12)(34), (13)(24), (14)(23) \} \leq S_4$.
        This is a normal subgroup as conjugation preserves cycle type
        and $K$ is the union of the conjugacy classes $(1)$ and $(2, 2)$.

        Now $|S_4| = 24$ and $|K| = 4$. By Lagrange's theorem, if $H$
        is a normal subgroup of $S_4$ containing $K$ its order must be divisible by 4
        and divide 24. Hence $|H| \in \{ 4, 8, 12, 24 \}$. Now since $H$ is normal it is the union
        of conjugacy classes in $S_4$. the conjugacy classes of $S_4$
        correspond to the cycle types $(1), (2), (2, 2), (3), (4)$ with sizes
        $1, 6, 3, 8, 6$. $K$ is already the union of $(1)$ and $(2, 2)$
        so by inspection it is clear that either $|H| = 4$ in which case $H = K$,
        or $|H| = 12$ in which case $H = A_4$ or $|H| = 24$ in which case $H = S_4$.
        
        If $H = K$ then $H/K \cong \{ e \}$.

        Else, if $H = A_4$ then $|A_4/K| = 3$ hence by Lagrange's theorem
        \[ A_4/K \cong \mathbb Z / 3\mathbb Z \]

        Finally if $H = S_4$ then we can use the action defined in lecture
        by $S_4 \curvearrowright \Delta$ where
        \[ \Delta = \{ P_1, P_2, P_3 \} \]
        with
        \begin{align*}
            P_1 &= \{ \{ 1, 2 \}, \{ 3, 4 \} \} \\
            P_2 &= \{ \{ 1, 3 \}, \{ 2, 4 \} \} \\
            P_3 &= \{ \{ 1, 4 \}, \{ 2, 3 \} \} \\
        \end{align*}
        and $S_4$ permutes the individual coordinates of each $P_i$.
        In lecture we showed that if $\pi: S_4 \to S_\Delta \cong S_3$
        is the permutation representation of this action, $\ker \pi = K$
        and so $|S_4/ K| = 6$ and since $|S_3| = 6$ this implies that $\pi$ is surjective
        and hence
        \[ S_4/K \cong S_3 \]

        Now to construct an explicit isomorphism
        \[ (S_4/K)/(A_4/K) \cong \{ -1, 1 \} \]
        we use question part (a) and map
        \[ \varphi: (S_4/K)/(A_4/K) \to S_4/A_4 \]
        given by
        \[ \varphi([[\sigma]_K]_{A_4/K}) = [\sigma]_{A_4} \]
        but since
        \[ S_4/A_4 \cong \{ -1, 1 \} \]
        via the sign map we can define an isomorphism
        \[ \psi: (S_4/K)/(A_4/K) \to \{ -1, 1 \} \]
        by
        \[ \psi([[\sigma]_K]_{A_4/K}) = \operatorname{sgn}(\sigma) \]
    \end{proof}

    \item[(c)]
    \begin{proof}
        Take $L = \{ e, (12)(34) \} \leq K$ which is normal in $K$
        but not normal in $G$ as conjugating by $(23)$ we get $(13)(24)$
        which is not in $L$. This means that $G/L$ is not a group
        and so the statement
        \[ (G/L)/(K/L) \cong G/K \]
        is not well defined as the set being quotiented is not a group.
    \end{proof}
\end{enumerate}

\end{document}